\section{অধ্যায়ভিত্তিক সৃজনশীল প্রশ্ন}
নিচের উদ্দীপকগুলো অধ্যায়ের ধারণা অনুশীলনের জন্য তৈরি।
\subsection{১. গ্রামের শিক্ষার্থী ও দূরের চিকিৎসক}
গ্রামের একজন শিক্ষার্থী ডাউনলোড করা পাঠ দেখে শেখে এবং ভিডিও সভায় শিক্ষককে প্রশ্ন করে। পরিবারের রোগীর ডিজিটাল প্রতিবেদন শহরের চিকিৎসককে পাঠানো হয়। অন্য একটি পরিবারের প্রয়োজনীয় যন্ত্র ও ব্যবহারের দক্ষতা নেই।
\begin{compactenum}
\item ক. বিশ্বগ্রাম কী?

\item খ. রেকর্ড করা পাঠ ও সরাসরি ক্লাস কীভাবে আলাদা?

\item গ. পরিবারের চিকিৎসায় ব্যবহৃত প্রযুক্তি ব্যাখ্যা কর।

\item ঘ. সেবাগুলো সবার সমান উপকার করছে কি?

\end{compactenum}

\subsection{২. দোকান ও স্বাধীন কর্মী}
একজন কারিগর নিজের দোকান থেকে ক্রেতাকে অনলাইনে পণ্য বিক্রি করেন। দোকানের হিসাবের কাজ চুক্তির ভিত্তিতে অন্য একটি প্রতিষ্ঠানকে দেন। তাঁর বোন বিভিন্ন ক্লায়েন্টের জন্য স্বাধীনভাবে গ্রাফিক নকশা করেন।
\begin{compactenum}
\item ক. E-commerce কী?

\item খ. Freelancing ও Outsourcing এক নয় কেন?

\item গ. দোকানের ব্যবসার ধরন ও হিসাব পরিচালনার ব্যবস্থাটি শনাক্ত কর।

\item ঘ. শুধু ওয়েবসাইট থাকলেই আয় বাড়বে কি?

\end{compactenum}

\subsection{৩. চালকের প্রশিক্ষণ}
প্রশিক্ষণকেন্দ্রে HMD পরে একজন শিক্ষার্থী কৃত্রিম রাস্তা দেখেন। স্টিয়ারিং ঘোরালে ভার্চুয়াল গাড়ি বাঁক নেয়; ট্র্যাকার মাথার অবস্থান মাপে। অনুশীলনে বাস্তব গাড়িকে দুর্ঘটনার ঝুঁকিতে ফেলতে হয় না।
\begin{compactenum}
\item ক. HMD-এর পূর্ণরূপ কী?

\item খ. Interaction কেন প্রয়োজন?

\item গ. প্রশিক্ষণব্যবস্থার কাজের ধাপ ব্যাখ্যা কর।

\item ঘ. VR বাস্তব প্রশিক্ষণের সম্পূর্ণ বিকল্প কি?

\end{compactenum}

\subsection{৪. কৃষি পরামর্শ ও ফসল গবেষণা}
একটি অ্যাপ পাতার ছবি বিশ্লেষণ করে সম্ভাব্য রোগের শ্রেণি দেখায়। গবেষক বিভিন্ন জাতের DNA-এর ক্রম তুলনা করে লবণসহিষ্ণুতার সঙ্গে সম্পর্কিত সম্ভাব্য জিন বাছাই করেন। পরে পরীক্ষাগারে DNA পরিবর্তনের পরীক্ষা করেন।
\begin{compactenum}
\item ক. Machine Learning কী?

\item খ. মডেলকে প্রশিক্ষণ দেওয়া ও প্রশিক্ষিত মডেলে ফল নির্ণয় করা আলাদা কেন?

\item গ. অ্যাপের কাজ ও DNA-এর ক্রম বিশ্লেষণের প্রযুক্তি শনাক্ত কর।

\item ঘ. কাঙ্ক্ষিত ফসল পেতে তিন কাজের সম্পর্ক মূল্যায়ন কর।

\end{compactenum}

\subsection{৫. বিপজ্জনক কারখানায় রোবট}
রোবটের ক্যামেরা গরম ধাতুর অবস্থান মাপে। নিয়ন্ত্রক নির্দেশ দেয়; মোটর বাহু চালায়; গ্রিপার ধাতু ধরে। কাজের ফল সেন্সর দিয়ে আবার মাপা হয়।
\begin{compactenum}
\item ক. Actuator কী?

\item খ. সেন্সর ও End Effector আলাদা কেন?

\item গ. উদ্দীপকের Sense--Think--Act ব্যাখ্যা কর।

\item ঘ. মানুষ পুরোপুরি বাদ দিলে কি কারখানা নিরাপদ?

\end{compactenum}

\subsection{৬. নির্বাচিত টিস্যুর চিকিৎসা}
চিকিৎসক ইমেজিংয়ের সাহায্যে অস্বাভাবিক টিস্যু শনাক্ত করেন। ক্রায়োপ্রোব দিয়ে নিয়ন্ত্রিত ঠান্ডা প্রয়োগ করেন। জমাট বাঁধানো ও গলানোর ধাপ পর্যবেক্ষণ করে আশপাশের টিস্যু রক্ষার চেষ্টা করেন।
\begin{compactenum}
\item ক. Cryogen কী?

\item খ. প্রোব ব্যবহার করলেই কোনো চিকিৎসাকে ক্রায়োসার্জারি বলা যায় না কেন?

\item গ. উদ্দীপকের চিকিৎসা ও ICT-এর ভূমিকা ব্যাখ্যা কর।

\item ঘ. পদ্ধতিটি সম্পূর্ণ ঝুঁকিমুক্ত কি?

\end{compactenum}

\subsection{৭. বন্যার মানচিত্র ও উদ্ধার}
উপগ্রহের সেন্সর বন্যার ছবি পাঠায়। GIS-এ রাস্তা ও আশ্রয়কেন্দ্রের স্তর মিলিয়ে পরিকল্পনা করা হয়। উদ্ধারকারী GPS-এ নিজের অবস্থান দেখেন। নিয়ন্ত্রণকেন্দ্র ড্রোনের ব্যাটারি ও তাপমাত্রার তথ্য গ্রহণ করে।
\begin{compactenum}
\item ক. Telemetry কী?

\item খ. GPS ও GIS আলাদা কেন?

\item গ. চারটি কাজের প্রযুক্তি শনাক্ত কর।

\item ঘ. তথ্য পেলেই সঠিক সিদ্ধান্ত নিশ্চিত কি?

\end{compactenum}

\subsection{৮. স্মার্ট উৎপাদন ও প্রতিরক্ষা}
কম্পিউটারে যন্ত্রাংশের নকশা করে CNC দিয়ে কাটা হয়। RFID দিয়ে মজুত পণ্য শনাক্ত করা হয়। অন্য একটি প্রতিষ্ঠানে রাডার লক্ষ্যবস্তু শনাক্ত করে এবং এনক্রিপ্ট করা যোগাযোগে তথ্য পাঠানো হয়।
\begin{compactenum}
\item ক. CAD-এর পূর্ণরূপ কী?

\item খ. RFID ও Barcode-এর Reader আলাদা কেন?

\item গ. উৎপাদনের ধাপগুলোতে ICT-এর কাজ ব্যাখ্যা কর।

\item ঘ. Radar ও Encryption কি একই নিরাপত্তা দেয়?

\end{compactenum}

\subsection{৯. পরিচয় যাচাই ও DNA গবেষণা}
কর্মী নিজের ID দেওয়ার পরে আঙুলের বৈশিষ্ট্য রেফারেন্স টেমপ্লেটের সঙ্গে মেলান। কারও আঙুলে ক্ষত থাকলে বিকল্প যাচাই লাগে। আলাদা পরীক্ষাগারে পারিবারিক সম্পর্ক যাচাইয়ের জন্য DNA প্রোফাইল তুলনা করা হয়; গবেষক BLAST ব্যবহার করেন।
\begin{compactenum}
\item ক. Biometric Template কী?

\item খ. প্রথম ব্যবস্থাটি Verification কেন?

\item গ. আঙুলের বৈশিষ্ট্য মেলানো, DNA প্রোফাইল তুলনা ও BLAST-এর কাজ ব্যাখ্যা কর।

\item ঘ. বায়োমেট্রিক্স পাসওয়ার্ডের চেয়ে নিরাপদ হলেই কি সব সমস্যা দূর হয়?

\end{compactenum}

\subsection{১০. ন্যানো বাহক ও দায়িত্বশীল প্রকাশ}
গবেষক ন্যানো মাত্রার বাহক দিয়ে নির্দিষ্ট টিস্যুতে ওষুধ পৌঁছানোর পরীক্ষা করেন। একজন শিক্ষার্থী প্রতিবেদনে উৎস উল্লেখ না করে অন্যের লেখা নিজের বলে দাবি করেন। আরেক ঘটনায় ব্যাংকের পরিচয় দিয়ে ভুয়া বার্তায় অ্যাকাউন্টের OTP চাওয়া হয়।
\begin{compactenum}
\item ক. এক ন্যানোমিটার কত মিটার?

\item খ. ছোট হলেই কোনো প্রযুক্তিকে ন্যানো প্রযুক্তি বলা যায় না কেন?

\item গ. তিনটি ঘটনার ধারণা শনাক্ত কর।

\item ঘ. বৈজ্ঞানিক অগ্রগতিতে নৈতিকতা কেন জরুরি?

\end{compactenum}
